%%
%% Chapter 4: Experiments
%% Independent file, included by ../main.tex
%%

\section{Experiments}
\label{sec:exp}

\subsection{Data and Protocol}
\label{sec:exp:protocol}

We evaluate on LUNA16 \cite{setio2017luna16} with the CBCT projection
simulator used in \cite{lin2024difgaussian}. The dataset is split into
738 training, 50 evaluation, and 100 test cases. Unless stated otherwise,
all models are trained from random initialization with the official
400-epoch protocol: 738 optimizer updates per epoch (batch size 1),
AdamW with $\mathrm{lr}=2{\times}10^{-4}$, per-epoch StepLR decay
$(10^{-3})^{1/400}$ per epoch, global gradient clipping at 1.0, FP32
precision, and a fixed random seed. The sparse context always consists
of the same six fixed projections with raw IDs $\{0,8,16,25,33,41\}$;
the dense training set contains 48 views. The network samples a random
sub-volume block and 1024 query points per step and is supervised in
volume space. Models are checkpointed every 20 epochs; the primary
comparison uses the final (epoch 400) model, and a supplementary
best-eval column reports the best PSNR over the 20 evaluation nodes.
Following the deterministic protocol, all probabilistic variants report
$N{=}8$ Monte-Carlo prior-mean estimates at inference. Metrics are peak
signal-to-noise ratio (PSNR, dB) and structural similarity (SSIM) on
the evaluation volumes at native scale.

\subsection{Implementation Details}
\label{sec:exp:impl}

All variants share the backbone of Section~\ref{sec:method:backbone}:
UNet encoder ($out\_ch=128$), geometry-transformer fusion with
parameters as in the released configuration (2 blocks, 4 heads, window
$\pm1$, feature scales $\{4,8,16\}$, $K{=}512$ anchors at resolution 12,
token initialization \texttt{max}). The probabilistic layer uses a
64-dimensional global latent; the FiLM modulation is initialized with
zero-mean weights of std $10^{-3}$ so that early training behaves like
the deterministic model. The GT encoder is a 3D CNN with channel
sizes $(1,16,16,32)$ and two stride-2 convolutions, followed by global
mean pooling, producing a 32-dimensional case descriptor. The posterior
head therefore receives a $128+7+32$-dimensional conditioning
vector. Prior and posterior heads share the bounded-Gaussian
parameterization with $\sigma\in[0.001,3.0]$ and
$\sigma_{\mathrm{init}}=0.5$. The KL is annealed linearly over 500 steps
to $\beta_f{=}2{\times}10^{-2}$ with a free-bits floor of $0.5$ nats;
the decoder conditions on the context set
(\texttt{posterior\_decoder\_views='context'}).

\subsection{Main Results}
\label{sec:exp:main}

Table~\ref{tab:results} reports reconstruction quality on the LUNA16
evaluation set. Several findings stand out.

\begin{table}[t]
\centering
\small
\setlength{\tabcolsep}{4pt}
\caption{Reconstruction quality on LUNA16 (six sparse views). Final
refers to the epoch-400 protocol model; best-eval reports the best PSNR
checkpoint over the 20 evaluation nodes. All probabilistic numbers use
the deployment-side estimator (prior MC mean, $N{=}8$), which never
touches dense views or ground truth.}
\label{tab:results}
\begin{tabular}{@{}p{3.3cm}p{2.5cm}ccc@{}}
\toprule
Model & Variant & PSNR (dB) & SSIM & Best-eval PSNR (dB) \\
\midrule
Deterministic backbone \cite{lin2024difgaussian} & -- & 29.54 & 0.870 & -- \\
\midrule
GT-CVAE & max-init, run~1 & 29.50 & 0.868 & 29.50 (E400) \\
GT-CVAE & max-init, run~2 & 29.06 & 0.854 & 29.06 (E400) \\
GT-CVAE & mean-init, run~1 & 27.55 & 0.808 & -- \\
GT-CVAE & mean-init, run~2 & 26.96 & 0.779 & -- \\
\bottomrule
\end{tabular}
\end{table}

\emph{The GT-conditioned CVAE reaches deterministic-level quality while
emitting uncertainty.}
The best GT-conditioned CVAE run reaches $29.50$\,dB / $0.868$\,SSIM at
the epoch-400 checkpoint, essentially matching the deterministic
backbone ($29.54$\,dB / $0.870$\,SSIM, a gap of only
$0.04$\,dB), while additionally emitting per-point predictive variance. Crucially, all of these numbers come from the
\emph{deployment} estimator ($z\sim p(z\mid C)$, no GT), so the reported
accuracy is what an operator would actually obtain.

\emph{The GT-conditioned posterior costs almost nothing.}
The model conditions its posterior on the ground-truth volume during
training, yet matches the deterministic backbone to within $0.04$\,dB
under the same \texttt{token\_init=max} setting. This confirms that
encoding the target into the variational posterior concentrates the
latent on case-specific content without overfitting the target, and
that the deployment-side prior retains the information needed for
accurate reconstruction.

\emph{Token initialization matters more than posterior design.} With
\texttt{token\_init\allowbreak=mean} both runs collapse to
$\approx 27$\,dB, a gap of roughly 2\,dB to the \texttt{max}-initialized
runs. All other settings are identical; the mechanism is analyzed in
Section~\ref{sec:exp:ablation}.

\subsection{Ablations}
\label{sec:exp:ablation}

\textbf{GT conditioning.}
Given the same \texttt{max} token initialization, the GT-conditioned
posterior (which encodes the ground truth during training) reaches
$29.06$--$29.50$\,dB against $29.54$\,dB for the deterministic
backbone, confirming that the GT-conditioned posterior provides the
intended information concentration without overfitting the target.

\textbf{Token initialization.}
Table~\ref{tab:results} shows a systematic effect:
\texttt{mean} \emph{vs.} \texttt{max} fusion of view tokens into
anchors changes the model's final PSNR by $\sim\!2$\,dB ($27.0$--$27.5$ vs.\
$29.1$--$29.5$\,dB). We attribute this to the interaction between token
initialization and the FiLM-gated latent: with \texttt{mean}
initialization the anchor features are more uniform, the fusion
transformer saturates earlier, and the latent has less discriminating
signal to modulate.